% =====================================================================
\chapter{Тренер и симуляторы}
\label{ch:trainer}
% =====================================================================
\chapterintro
  {как соединить два пульта в режиме «учитель--ученик», как использовать пульт
   в компьютерном симуляторе через USB или Bluetooth, как пользоваться EdgeTX Companion}
  {второй пульт и тренерский кабель (для тренера); компьютер (для симулятора)}

\section{Режим «учитель--ученик»}

Учитель держит пульт, связанный с моделью, и передаёт управление ученику кнопкой.
Отпустил кнопку --- управление снова у учителя. Это самый безопасный способ научиться летать
на настоящей модели.

\begin{figure}[H]
\centering
\begin{tikzpicture}[node distance=10mm and 18mm]
  \node[blk,fill=blkin,minimum width=26mm] (s) {Пульт ученика\\\scriptsize Slave};
  \node[blk,fill=blkmix,right=of s,minimum width=26mm] (m) {Пульт учителя\\\scriptsize Master};
  \node[blk,fill=blkout,right=of m,minimum width=20mm] (r) {Модель};
  \draw[arr] (s) -- node[above,lbl]{кабель/радио} node[below,lbl]{PPM, каналы} (m);
  \draw[arr,decorate,decoration={snake,amplitude=1.2pt,segment length=5pt,post length=3pt}] (m) -- (r);
  \node[below=4mm of m,lbl,text width=40mm,align=center] {кнопка «Trainer»:\\нажата --- управляет ученик};
\end{tikzpicture}
\caption{Схема тренерского режима}
\end{figure}

\subsection{Подключение}

Варианты связи пультов:
\begin{itemize}
  \item \textbf{Кабель} 3,5\,мм «джек--джек» через гнёзда тренера (если на вашей ревизии
        пульта есть гнездо тренера; уточните по инструкции производителя).
        Режимы \menu{Master/Jack} и \menu{Slave/Jack}.
  \item \textbf{Мультипротокольный модуль} учителя может \emph{принимать} сигнал пульта ученика
        по радио (режим \menu{Master/Multi}); ученик при этом работает как обычный передатчик
        FrSky и т.\,п.
  \item \textbf{Последовательный порт} (\menu{Master/Serial}) --- с внешними приёмниками
        или модулями беспроводного тренера (например, на основе ELRS Backpack).
\end{itemize}

\subsection{Пульт ученика}

\begin{enumerate}
  \item Создайте отдельную модель \code{Trainee}.
  \item \mpath{MDL\sep Setup}: радиомодули выключить (\menu{Internal RF} = \menu{OFF}),
        \menu{Trainer mode} = \menu{Slave/Jack}.
  \item Порядок каналов --- такой же, как у учителя (AETR). Никаких миксеров для крыла и т.\,п. ---
        только «чистые» стики: миксует пульт учителя.
\end{enumerate}

\subsection{Пульт учителя}

\begin{enumerate}
  \item В модели: \mpath{MDL\sep Setup\sep Trainer mode} = \menu{Master/Jack}.
  \item \mpath{SYS\sep Trainer}: для каждой оси (\sw{Ail}, \sw{Ele}, \sw{Thr}, \sw{Rud}) режим
        \menu{:=} (заменять стик учителя) или \menu{+=} (добавлять к стику учителя), вес
        (например, 100\,\% или меньше для начинающего) и номер канала ученика.
  \item Там же --- калибровка: поставьте стики ученика в центр и выберите \menu{Cal}.
  \item Специальная функция: \menu{Switch} = кнопка без фиксации (\sw{SF} на Boxer),
        \menu{Function} = \menu{Trainer}, \menu{Sticks}.
\end{enumerate}

\begin{caution}
Перед полётом проверьте на земле: при нажатой кнопке рули и газ должны слушаться
пульта ученика, при отпущенной --- пульта учителя. Проверьте и газ, и направление всех осей.
\end{caution}

\section{Пульт в симуляторе}

Тренироваться в симуляторе --- дёшево и эффективно: разбитый виртуальный квадрокоптер
ничего не стоит.

\subsection{USB-джойстик}

\begin{enumerate}
  \item Создайте модель \code{Sim}: радиомодуль \menu{OFF}, каналы \sw{CH1}--\sw{CH4} ---
        стики, \sw{CH5}--\sw{CH8} --- переключатели (для сброса, режимов и т.\,п.).
  \item Подключите пульт к компьютеру по USB и выберите режим \menu{USB Joystick (HID)}.
  \item Компьютер увидит игровой контроллер с 8 осями (и кнопками --- в расширенном режиме).
  \item В симуляторе откалибруйте оси и назначьте функции.
\end{enumerate}

В \mpath{MDL\sep Setup\sep USB Joystick} можно выбрать \menu{Classic} (каналы → оси) или
\menu{Advanced} --- тогда для каждого канала задаётся, будет он осью, кнопкой или
переключателем, что удобно для игр и симуляторов с кнопками.

\subsection{Беспроводной джойстик ELRS}

Модуль ExpressLRS на ESP32 умеет работать как Bluetooth-геймпад. Режим включается в скрипте
\mpath{SYS\sep Tools\sep ExpressLRS}, пункт \menu{BLE Joystick}. Удобно для симуляторов на планшете или ноутбуке без проводов.
Пока режим активен, модуль не управляет моделью.

\subsection{Какие симуляторы бывают}

\begin{itemize}
  \item FPV-коптеры: Liftoff, VelociDrone, Uncrashed, TRYP FPV, DRL Simulator и др.;
  \item самолёты, планеры, вертолёты: RealFlight, PicaSim, ClearView, Heli-X и др.
\end{itemize}

\begin{tip}
В FPV-симуляторе задайте такие же рейты (rates), как в Betaflight на настоящем
коптере. Тогда навыки переносятся один в один.
\end{tip}

\section{EdgeTX Companion}

\textbf{Companion} --- программа для компьютера (Windows, macOS, Linux):
\begin{itemize}
  \item чтение и запись моделей и настроек пульта через SD-карту;
  \item редактирование моделей в удобном оконном интерфейсе;
  \item \textbf{симулятор пульта}: можно проверить миксеры, логические переключатели и
        Lua-скрипты, не включая настоящий пульт;
  \item печать и сравнение моделей, резервные копии.
\end{itemize}

Для учебной группы Companion особенно удобен: преподаватель готовит «эталонную» модель,
а ученики загружают её себе.

\begin{summary}
\begin{itemize}
  \item Тренер: ученик --- \menu{Slave}, учитель --- \menu{Master} + функция \menu{Trainer}
        на кнопке без фиксации.
  \item Симулятор: отдельная модель, USB-джойстик или Bluetooth через ELRS.
  \item Companion --- редактор и симулятор EdgeTX на компьютере.
\end{itemize}
\end{summary}
